% Revize 2026-10-27: cílené doplnění; nové úlohy a řešení jsou vlastní tvorba.
% Metodický vzor: 4_rocnik/pripravy/Uvod_do_kombinatoriky.pdf (samostatný dokument).
% Hodina 34, 1. ročník SPŠ. Rozsah: Rozpis_vyuky_MAT_1_rocnik.tex, řádek 34.

% Plán: 1_rocnik/tematicke_plany/SPS.pdf; společná projektová šablona.

% Učebnice: ucebnice/01_Zakladni_poznatky/Mocniny_a_odmocniny/Mocniny a Odmocniny.tex

% Původ úloh: P1: učebnice, „Částečné odmocňování“, oba příklady. Ostatní vlastní úlohy.

% Metodický zdroj: externi_zdroje/realisticky/Odmocniny/07_Druha_odmocnina.pdf

% URL: http://www.realisticky.cz/ucebnice/01%20Matematika%20S%C5%A0/01%20Z%C3%A1kladn%C3%AD%20poznatky/02%20%C4%8C%C3%ADseln%C3%A9%20obory/07%20Druh%C3%A1%20odmocnina.pdf

% Metodický zdroj: externi_zdroje/realisticky/Odmocniny/08_Treti_odmocnina.pdf

% URL: http://www.realisticky.cz/ucebnice/01%20Matematika%20S%C5%A0/01%20Z%C3%A1kladn%C3%AD%20poznatky/02%20%C4%8C%C3%ADseln%C3%A9%20obory/08%20T%C5%99et%C3%AD%20odmocnina.pdf

\documentclass[11pt,a4paper]{article}

% Spolecny styl pro ucitelske pripravy MAT (Work/MAT)
% Verze 1.1, 2026-10-23
\NeedsTeXFormat{LaTeX2e}
\ProvidesPackage{priprava}[2026/10/23 v1.1 Sablona ucitelskych priprav MAT]

% Zakladni typografie a cestina
\RequirePackage[T1]{fontenc}
\RequirePackage[utf8]{inputenc}
\RequirePackage[czech]{babel}
\RequirePackage{lmodern}
\RequirePackage{microtype}
\RequirePackage{amsmath,amssymb}
\RequirePackage{xcolor}
\RequirePackage{enumitem}
\RequirePackage[most]{tcolorbox}
\RequirePackage{tabularx,booktabs,array}
\RequirePackage{fancyhdr}
\RequirePackage{needspace}
\RequirePackage{geometry}
\RequirePackage[hidelinks]{hyperref}

\geometry{a4paper,top=17mm,bottom=18mm,left=23mm,right=23mm}
\definecolor{notegray}{gray}{0.39}
\setlength{\parindent}{0pt}
\setlength{\parskip}{2.5pt plus 1pt minus 0.5pt}
\setlist[itemize]{topsep=2pt,itemsep=1pt,parsep=0pt,leftmargin=1.45em}
\setlist[enumerate]{topsep=2pt,itemsep=2pt,parsep=0pt,leftmargin=1.65em}
\setcounter{secnumdepth}{0}
\renewcommand{\arraystretch}{1.12}

% Zahlavi: vlevo rocnik | globalni tema, vpravo cislo: nazev hodiny
% Zapati zustava zcela prazdne.
\pagestyle{fancy}
\fancyhf{}
\renewcommand{\headrulewidth}{0.35pt}
\renewcommand{\footrulewidth}{0pt}
\setlength{\headheight}{14pt}
\newcommand{\nastavPripravu}[4]{%
  \fancyhead[L]{\small #1\enspace\textbar\enspace #2}%
  \fancyhead[R]{\small #3: #4}%
}

% Minimalisticke cernobile boxy. Kazda definice zvlast.
\tcbset{commonlesson/.style={enhanced,breakable=false,colback=white,colframe=black,
  colbacktitle=white,coltitle=black,boxrule=0.45pt,
  arc=0.8mm,outer arc=0.8mm,left=3mm,right=3mm,top=1.6mm,bottom=1.6mm,
  toptitle=1mm,bottomtitle=0.7mm,fonttitle=\bfseries,
  attach boxed title to top left={xshift=2.5mm,yshift=-1.15mm},
  boxed title style={colback=white,colframe=white,boxrule=0pt,left=1mm,right=1mm}}}
\newtcolorbox{zadaniBox}[1][]{commonlesson,title={Zadání},#1}
\newtcolorbox{prikladBox}[1][]{commonlesson,title={Příklad},#1}
\newtcolorbox{definiceBox}[1][]{commonlesson,title={Definice},#1}

% Spolecny box pro zadani a jeho reseni, bez mezery mezi castmi.
% Pouziti:
%   \begin{prikladReseniBox}
%     Zadani ulohy.
%     \tcblower
%     \begin{align*} ... \end{align*}
%   \end{prikladReseniBox}
% Spodni cast ma jemny sedy podklad a delici linku; lze v ni sazet
% vice radku matematiky. Obe casti maji jen jeden spolecny ramecek.
\tcbset{solutionlesson/.style={
  commonlesson,
  enhanced jigsaw,
  breakable,
  colbacklower=black!3,
  segmentation style={solid,draw=black!25,line width=0.35pt},
  fontlower=\small\color{notegray},
  before lower={\noindent\textbf{Řešení:}\enspace},
  leftlower=3mm,rightlower=3mm
}}
\newtcolorbox{prikladReseniBox}[1][]{solutionlesson,title={Příklad},#1}
\newtcolorbox{zadaniReseniBox}[1][]{solutionlesson,title={Zadání},#1}

% Pomucky pro autora pripravy: sedive poznamky; \reseni pro starsi pripravy
\newcommand{\poznamka}[1]{\par\begingroup\small\color{notegray}\itshape
  \textnormal{\textbf{Poznámka pro učitele.}}\enspace #1\par\endgroup}
\newcommand{\reseni}[1]{\par\begingroup\small\color{notegray}
  \textbf{Očekávané řešení:}\enspace #1\par\endgroup}
\newcommand{\blok}[1]{\needspace{5\baselineskip}\subsection*{#1}}
\newcommand{\vyrok}[1]{\uv{#1}}


\nastavPripravu{1. ročník}{Odmocniny}{34}{Částečné odmocňování}

\begin{document}

\blok{Výklad a souvislosti}

Odmocnina z čísla, které není úplnou druhou mocninou, může obsahovat zbytečně velký radikand. Vyhledej v něm čtverec a použij pravidlo ze 33. hodiny. Až po této úpravě hledej stejné odmocniny ve více členech.

\begin{definiceBox}[title={Definice},breakable=false]
\textbf{Částečné odmocnění} je úprava, při níž z radikandu vyjmeme úplnou mocninu odpovídající odmocniteli a její odmocninu zapíšeme před odmocnítko.
\end{definiceBox}


Pro reálné $a$ a $b\ge0$ platí $\sqrt{a^2b}=|a|\sqrt b$. U třetí odmocniny je $\sqrt[3]{a^3b}=a\sqrt[3]b$ pro všechna reálná $a,b$. Členy se stejnou odmocninou slučujeme pomocí vytýkání: $u\sqrt b+v\sqrt b=(u+v)\sqrt b$, což je pouze distributivní zákon.

\blok{Řešené příklady}

\begin{prikladReseniBox}[title={Příklad 1},breakable=false]
Částečně odmocni $\sqrt{72}$ a $\sqrt[3]{54}$.
\tcblower
$\sqrt{72}=\sqrt{36\cdot2}=6\sqrt2$; $\sqrt[3]{54}=\sqrt[3]{27\cdot2}=3\sqrt[3]2$.
\end{prikladReseniBox}


\begin{prikladReseniBox}[title={Příklad 2},breakable=false]
Uprav $\sqrt{48}-2\sqrt{12}+\sqrt{75}$ a $\sqrt{50x^2}$.
\tcblower
První výraz je $4\sqrt3-4\sqrt3+5\sqrt3=5\sqrt3$. Druhý je pro všechna reálná $x$ roven $5|x|\sqrt2$, protože $\sqrt{x^2}=|x|$.
\end{prikladReseniBox}


\blok{Samostatné procvičení}

\begin{zadaniReseniBox}[title={Procvičení},breakable=false]
\begin{enumerate}[label=\alph*)]
\item Uprav $\sqrt{98}$ a $\sqrt[3]{-128}$.
\item Uprav $3\sqrt8-\sqrt{18}+2\sqrt2$.
\item Uprav $\sqrt{18a^2}$ pro reálné $a$.
\end{enumerate}
\tcblower
\begin{enumerate}[label=\alph*)]
\item $7\sqrt2$ a $-4\sqrt[3]2$; použijeme $98=49\cdot2$ a $-128=(-4)^3\cdot2$.
\item $6\sqrt2-3\sqrt2+2\sqrt2=5\sqrt2$.
\item $3|a|\sqrt2$. Absolutní hodnotu nelze vynechat například pro $a=-1$.
\end{enumerate}
\end{zadaniReseniBox}


\poznamka{Nejprve nech žáky hledat největší čtvercový dělitel, potom slučovat členy. Různé odmocniny nesčítej jako radikandy: $\sqrt2+\sqrt3$ není $\sqrt5$. Při pomalejším tempu ponech třetí odmocninu v procvičení jako rezervu.}


\blok{Navazující procvičování ve třídě}

\begin{zadaniReseniBox}[title={Základ: najdi vhodnou mocninu pod odmocnítkem},breakable=false]
\begin{enumerate}[label=\alph*)]
\item Částečně odmocni $\sqrt{108}$, $\sqrt{300}$, $\sqrt[3]{250}$ a $\sqrt[3]{-432}$.
\item Uprav $2\sqrt{27}+\sqrt{48}-\sqrt{75}$ a $\sqrt{32}+\sqrt{18}$.
\item Uprav $\sqrt{72x^2}$ a $\sqrt[3]{16a^3}$ pro reálné proměnné.
\end{enumerate}
\tcblower
\begin{enumerate}[label=\alph*)]
\item $6\sqrt3$, $10\sqrt3$, $5\sqrt[3]2$, $-6\sqrt[3]2$. Rozklady využívají $36$, $100$, $125$ a $-216$.
\item $6\sqrt3+4\sqrt3-5\sqrt3=5\sqrt3$; $4\sqrt2+3\sqrt2=7\sqrt2$.
\item $6|x|\sqrt2$ a $2a\sqrt[3]2$. U liché odmocniny absolutní hodnota nepatří.
\end{enumerate}
\end{zadaniReseniBox}

\begin{zadaniReseniBox}[title={Rezerva: stejné odmocniny až po úpravě},breakable=false]
\begin{enumerate}[label=\alph*)]
\item Uprav $\sqrt{200}-3\sqrt8+2\sqrt{50}$ a $\sqrt{a^4}$ pro reálné $a$.
\item Uprav $\sqrt{12}\sqrt{27}-\sqrt{75}/\sqrt3$.
\end{enumerate}
\tcblower
\begin{enumerate}[label=\alph*)]
\item $10\sqrt2-6\sqrt2+10\sqrt2=14\sqrt2$; $\sqrt{(a^2)^2}=|a^2|=a^2$.
\item $\sqrt{324}-\sqrt{25}=18-5=13$.
\end{enumerate}
\end{zadaniReseniBox}


\poznamka{Základ tvoří původní ukázky a nová první skupina; u každé odmocniny nech nejprve navrhnout rozklad. Další dva složené výpočty jsou rezerva. Domácí pětice kombinuje rozklad, slučování a znaménko proměnné. Nezaměňuj koeficient před odmocninou s číslem pod odmocnítkem.}

\blok{Domácí procvičení}

\begin{zadaniReseniBox}[title={Domácí úlohy},breakable=false]
\begin{enumerate}[label=\alph*)]
\item Částečně odmocni $\sqrt{147}$ a $\sqrt[3]{-375}$.
\item Uprav $3\sqrt{12}-\sqrt{27}+\sqrt{75}$.
\item Uprav $\sqrt{8x^2}$ a ověř pro $x=-2$.
\item Uprav $\sqrt[3]{54a^3}$ pro reálné $a$.
\item Uprav $2\sqrt{45}-\sqrt{20}+\sqrt{80}$.
\end{enumerate}
\tcblower
\begin{enumerate}[label=\alph*)]
\item $7\sqrt3$ a $-5\sqrt[3]3$.
\item $6\sqrt3-3\sqrt3+5\sqrt3=8\sqrt3$.
\item $2|x|\sqrt2$; pro $x=-2$ je na obou stranách $4\sqrt2$.
\item $3a\sqrt[3]2$.
\item $6\sqrt5-2\sqrt5+4\sqrt5=8\sqrt5$.
\end{enumerate}
\end{zadaniReseniBox}

\end{document}
